%\thispagestyle{fancy}
\pagecolor{gray!7}
\chapter{常见函数的拉氏变换}

\begin{table}[H]
	\centering
	\caption{常见函数的拉氏变换速查表}
	\label{tab:changjianhanshu-laplas}
	% 关键：缩放行高为原来的1.5倍（1.3~1.8之间都合适，越大边距越大）
	\renewcommand{\arraystretch}{1.25}
	\resizebox{0.4\textwidth}{!}{%
		\begin{tabular}{c|c}
			\hline \hline
		  像函数 $F\left(s\right)即L\left(y\right)$	& 原函数$f\left(x\right)$   \\ \hline
		  1	& $\delta\left(x\right)$ \\ \hline
		  $\dfrac{1}{s}$ & 1   \\ \hline
		  $\dfrac{1}{s^{2}}$ & $x$ \\ \hline
		  $\dfrac{n!}{s^{n+1}}$ & $x^{n}$  \\ \hline
		  $\dfrac{1}{s-a}$ & $e^{ax}$ \\ \hline
		  $\dfrac{n!}{\left(s-a\right)^{n+1}}$  &  $x^{n}\cdot e^{ax}$     \\ \hline
		  $\dfrac{\omega}{s^{2}+\omega^{2}}$ & $sin\left(\omega x\right)\:$   \\ \hline
		  $\dfrac{s}{s^{2}+\omega^{2}}$ & $cos\left(\omega x\right)$  \\ \hline
		  $\dfrac{s^{2}-\omega^{2}}{\left(s^{2}+\omega^{2}\right)^{2}}$ & $x\cdot cos\left(\omega x\right)$  \\ \hline
		  $\dfrac{2s\omega}{\left(s^{2}+\omega^{2}\right)^{2}}$ & $x\cdot sin\left(\omega x\right)$ 
		   \\ \hline \hline
		\end{tabular}%
	}
\end{table}

\vspace{-10pt}

\begin{spacing}{2.2}
	表格\ref{tab:changjianhanshu-laplas}中总结了常见函数的拉式变换。尤其注意对$e^{ax}$的拉式变换，可看作是$e^{ax}\cdot1$，得结果为$\dfrac{1}{s-a}$，可以看作是函数$e^{ax}$对函数1作用后，函数1的拉式变换像函数由原来的$\dfrac{1}{s}$,变为了$\dfrac{1}{s-a}$，也即平移了$a$个单位。类似的，原函数$x^{n}$经拉氏变换后的像函数为$\dfrac{n!}{s^{\left(n+1\right)}}$，经过函数$e^{ax}$对函数$x^{n}$作用后，$x^{n}$的像函数变为$\dfrac{n!}{\left(s-a\right)^{n+1}}$，也即平移了$a$个单位。此处函数$e^{ax}$对函数$x、x^{2}、x^{n}、sin\left(\omega x\right)、cos\left(\omega x\right)$等进行作用，在其像函数上体现为平移$a$个单位，这样的性质称为"位移性质"或"位移定理"，之后的"位移性质"章节将进一步举例说明
\end{spacing}

\section{对$x^{n}$的拉氏变换}

\begin{spacing}{2.2}
	\begin{center}
		$L\left(x^{n}\right)=\int_{0}^{\infty}x^{n}\cdot e^{-sx}dx=\dfrac{n!}{s^{\left(n+1\right)}}$
		
		特别地，$n=0$时，$L\left(1\right)=\dfrac{1}{s}$ 
		
		特别地，$s=1$时，$\int_{0}^{\infty}x^{n}\cdot e^{-x}dx$，即为$\gamma$函数
		
		关于类$\gamma$函数，还需记住两个公式(通常用于概率论中)
		
		$\int_{-\infty}^{+\infty}e^{-\alpha x^{2}}dx=2\int_{0}^{+\infty}e^{-\alpha x^{2}}dx=\sqrt{\dfrac{\pi}{\alpha}}$  
		
		此处$\alpha=1$时即为经典的概率论公式
		
		$\int_{-\infty}^{+\infty}x^{2}\cdot e^{-\alpha x^{2}}dx=2\int_{0}^{+\infty}x^{2}\cdot e^{-\alpha x^{2}}dx=\dfrac{1}{2\alpha}\sqrt{\dfrac{\pi}{\alpha}}$
		
		至于 \hspace{1em} $\int_{0}^{+\infty}x\cdot e^{-\alpha x^{2}}dx$ ，\hspace{0.5em} 凑微分
		
		\hspace{3em} $\int_{0}^{+\infty}x^{3}\cdot e^{-\alpha x^{2}}dx$ ，\hspace{0.5em} 凑微分
		
		此式适用于气体在平衡态下的麦克斯韦-玻尔兹曼速率分布定律，推导平均速率
		
		\hspace{3em} $\int_{0}^{+\infty}x^{4}\cdot e^{-\alpha x^{2}}dx$ ，\hspace{0.5em} 不必掌握
	\end{center}
\end{spacing}

















